% =====================================================================
%  Bibliographie et netographie
% =====================================================================
%  Tenue a la main plutot qu'avec biblatex : le guide impose un format
%  precis (auteurs, titre, editeur, date, URL et date de consultation) et
%  la compilation reste un simple pdflatex lance deux fois.
% =====================================================================
\cleardoublepage
\phantomsection
\addcontentsline{toc}{chapter}{Bibliographie et netographie}
\chapter*{Bibliographie et netographie}
\markboth{Bibliographie et netographie}{}

\section*{Ouvrages}
\addcontentsline{toc}{section}{Ouvrages}

\begin{enumerate}[label={[\arabic*]},leftmargin=*,itemsep=4pt]

\item MARTIN Robert C., \textit{Clean Architecture: A Craftsman's Guide
  to Software Structure and Design}, Prentice Hall, 2017, 432 p.

\item NEWMAN Sam, \textit{Building Microservices: Designing Fine-Grained
  Systems}, 2\textsuperscript{e} edition, O'Reilly Media, 2021, 616 p.

\item EVANS Eric, \textit{Domain-Driven Design: Tackling Complexity in
  the Heart of Software}, Addison-Wesley, 2003, 560 p.

\item ROQUES Pascal, \textit{UML 2 par la pratique : etudes de cas et
  exercices corriges}, 7\textsuperscript{e} edition, Eyrolles, 2009,
  396 p.

\item AUBRY Claude, \textit{Scrum : le guide pratique de la methode agile
  la plus populaire}, 5\textsuperscript{e} edition, Dunod, 2019, 336 p.

\end{enumerate}

\section*{Articles et specifications}
\addcontentsline{toc}{section}{Articles et specifications}

\begin{enumerate}[label={[\arabic*]},leftmargin=*,itemsep=4pt,start=6]

\item CHEN Tianqi et GUESTRIN Carlos, \guillemotleft{} XGBoost: A Scalable
  Tree Boosting System \guillemotright{}, \textit{Proceedings of the
  22\textsuperscript{nd} ACM SIGKDD International Conference on Knowledge
  Discovery and Data Mining}, San Francisco, 2016, p. 785-794.

\item LEWIS Patrick \textit{et al.}, \guillemotleft{} Retrieval-Augmented
  Generation for Knowledge-Intensive NLP Tasks \guillemotright{},
  \textit{Advances in Neural Information Processing Systems}, vol. 33,
  2020, p. 9459-9474.

\item OBJECT MANAGEMENT GROUP, \textit{Unified Modeling Language
  Specification, version 2.5.1}, decembre 2017.

\item SCHWABER Ken et SUTHERLAND Jeff, \textit{The Scrum Guide}, novembre
  2020, 13 p.

\end{enumerate}

\section*{Netographie}
\addcontentsline{toc}{section}{Netographie}

\begin{enumerate}[label={[\arabic*]},leftmargin=*,itemsep=4pt,start=10]

\item MICROSOFT, \textit{Documentation .NET 8 et ASP.NET Core},
  \url{https://learn.microsoft.com/dotnet/}, consulte en mars 2026.

\item MICROSOFT, \textit{Entity Framework Core},
  \url{https://learn.microsoft.com/ef/core/}, consulte en mars 2026.

\item MASSTRANSIT, \textit{Documentation officielle},
  \url{https://masstransit.io/documentation}, consulte en avril 2026.

\item META, \textit{Documentation React},
  \url{https://react.dev/}, consulte en avril 2026.

\item VERCEL, \textit{Documentation Next.js},
  \url{https://nextjs.org/docs}, consulte en mai 2026.

\item EXPO, \textit{Documentation Expo et React Native},
  \url{https://docs.expo.dev/}, consulte en mai 2026.

\item POSTGRESQL GLOBAL DEVELOPMENT GROUP, \textit{PostgreSQL 16
  Documentation}, \url{https://www.postgresql.org/docs/16/}, consulte en
  mars 2026.

\item BROADCOM, \textit{RabbitMQ Documentation},
  \url{https://www.rabbitmq.com/docs}, consulte en avril 2026.

\item QDRANT, \textit{Qdrant Documentation},
  \url{https://qdrant.tech/documentation/}, consulte en juillet 2026.

\item DOCKER INC., \textit{Docker Compose Documentation},
  \url{https://docs.docker.com/compose/}, consulte en mars 2026.

\item PLANTUML, \textit{Reference Guide},
  \url{https://plantuml.com/guide}, consulte en juin 2026.

\end{enumerate}
